% Part: first-order-logic
% Chapter: syntax-and-semantics
% Section: free-vars-sentences

\documentclass[../../../include/open-logic-section]{subfiles}

\begin{document}

\olfileid{fol}{syn}{fvs}

\olsection{ازاد \printtoken{P}{variable} او \printtoken{P}{sentence}}

\begin{defn}[د !!a{variable} ازاد موردونه]
\ollabel{defn:free-occ}
په !!a{formula} کښې د !!a{variable} \emph{ازاد} موردونه په استقرايي
ډول داسې تعريفېږي:
\begin{enumerate}
\item \indcase*{!A}{$!A$ اټومي دے}{په $\indfrm$ کښې د
  !!{variable} ټول موردونه ازاد دي۔}

\tagitem{prvNot}{\indcase{!A}{\lnot !B}{د $\indfrm$ د !!{variable}
  ازاد موردونه کټ مټ د $!B$ ازاد موردونه دي۔}}{}

\item \indcase{!A}{(!B \ast !C)}{د $\indfrm$ د !!{variable} ازاد
  موردونه په $!B$ کښې ازاد موردونه او ورسره په~$!C$ کښې ازاد موردونه
  دي۔}

\tagitem{prvAll}{\indcase{!A}{\lforall[x][!B]}{په $\indfrm$ کښې د
  !!{variable} ازاد موردونه په~$!B$ کښې ټول ازاد موردونه دي، د~$x$
  له موردونو پرته۔}}{}

\tagitem{prvEx}{\indcase{!A}{\lexists[x][!B]}{په $\indfrm$ کښې د
  !!{variable} ازاد موردونه په~$!B$ کښې ټول ازاد موردونه دي، د~$x$
  له موردونو پرته۔}}{}
\end{enumerate}
\end{defn}

\begin{defn}[تړلي متغيرونه]
په يوه فارمول~$!A$ کښې د !!a{variable} يو مورد هله \emph{تړلے} دے
چې ازاد نۀ وي۔
\end{defn}

\begin{prob}
په \olref[fol][syn][fvs]{defn:free-occ} پسې د تړلو متغيرونو د موردونو
يو استقرايي تعريف ورکړئ۔
\end{prob}

\begin{defn}[ساحه]
\iftag{prvAll}{که $\lforall[x][!B]$ په يوه فارمول~$!A$ کښې د يوه
  فرعي فارمول مورد وي، نو په~$!A$ کښې د~$!B$ اړوند مورد د
  $\lforall[x]$ د اړوند مورد \emph{ساحه} بلل کېږي۔
  \iftag{prvEx}{د $\lexists[x]$ دپاره هم همدا شان۔}{}}{که
  $\lexists[x][!B]$ په يوه فارمول~$!A$ کښې د يوه فرعي فارمول مورد
  وي، نو په~$!A$ کښې د~$!B$ اړوند مورد د $\lexists[x]$ د اړوند مورد
  \emph{ساحه} بلل کېږي۔}

که $!B$ په~$!A$ کښې د سور د يوه مورد
\iftag{prvAll}{$\lforall[x]$\iftag{prvEx}{ يا
    $\lexists[x]$}{}}{$\lexists[x]$} ساحه وي، نو په~$!B$ کښې د
$x$ ازاد موردونه په \iftag{prvAll}{$\lforall[x][!B]$\iftag{prvEx}{ او
$\lexists[x][!B]$}{}}{$\lexists[x][!B]$} کښې تړلي دي۔ وايو چې دا
موردونه د ياد شوي سور د مورد \emph{له خوا تړل شوي} دي۔
\end{defn}

\begin{ex}
لاندې فارمول په پام کښې ونيسئ:
\[
\lexists[\Obj v_0][\underbrace{\Atom{\Obj A^2_0}{\Obj v_0,\Obj v_1}}_{!B}]
\]
$!B$ د $\lexists[\Obj v_0]$ ساحه ده۔ سور په $!B$ کښې د $\Obj v_0$
مورد تړي، خو د $\Obj v_1$ مورد نۀ تړي۔ نو په دې حالت کښې $\Obj v_1$
ازاد متغير دے۔


اوس وينو چې په يوه زيات پېچلي !!{formula}~$!A$ کښې دا څنګه کار کوي:
\[
\lforall[\Obj v_0][\underbrace{(\Atom{\Obj A^1_0}{\Obj v_0} \lif
    \Atom{\Obj A^2_0}{\Obj v_0, \Obj v_1})}_{!B}] \lif \lexists[\Obj
  v_1][\underbrace{(\Atom{\Obj A^2_1}{\Obj v_0, \Obj v_1} \lor \lforall[\Obj v_0][\overbrace{\lnot \Atom{\Obj A^1_1}{\Obj v_0}}^{!D}])}_{!C}]
\]
$!B$ د لومړي $\lforall[\Obj v_0]$ ساحه، $!C$ د
$\lexists[\Obj v_1]$ ساحه، او $!D$ د دويم $\lforall[\Obj v_0]$ ساحه
ده۔ لومړے $\lforall[\Obj v_0]$ په~$!B$ کښې د $\Obj v_0$ موردونه
تړي، $\lexists[\Obj v_1]$ په $!C$ کښې د $\Obj v_1$ مورد تړي، او دويم
$\lforall[\Obj v_0]$ په~$!D$ کښې د $\Obj v_0$ مورد تړي۔ د $\Obj v_1$
لومړے مورد او د $\Obj v_0$ څلورم مورد په~$!A$ کښې ازاد دي۔ د
$\Obj v_0$ وروستے مورد په $!D$ کښې ازاد، خو په $!C$ او~$!A$ کښې
تړلے دے۔
\end{ex}

\begin{defn}[جمله]
!!^a{formula}~$!A$ هله او يوازې هله يوه \article{sentence}
\emph{!!{sentence}} ده چې د !!{variable}s هېڅ ازاد مورد ونۀ لري۔
\end{defn}

% add examples!

\end{document}
